\documentclass{fisatprojectfinal}
\renewcommand{\baselinestretch}{1.25}
\usepackage{setspace}
\usepackage[utf8]{inputenc}
\usepackage{xcolor}
\usepackage{amssymb}

\title{Intervue -- AI-Powered Mock Interview and Career Readiness Platform}
\team{Pagau S \\ Gautham Krishna \\ Hussain \\ Safwan}
\author{Pagau S}
\regno{FIT23CSD065}
\bibliographystyle{plain}

\begin{document}
\maketitle
\makecert

\newpage
\pagenumbering{roman}
\setcounter{page}{1}
\newgeometry{top=4cm,bottom=0.1cm}
\thispagestyle{empty}
\renewcommand\abstractname{ABSTRACT}

% ---------------------------------------------------------------
% FIX: single \begin{abstract} — no nesting
% ---------------------------------------------------------------
\begin{abstract}
\vspace{2cm}
\begin{onehalfspace}
Intervue is an AI-powered mock interview and career readiness platform designed to address
the growing skills-confidence gap among fresh engineering graduates preparing for technical
job interviews. Many candidates enter the job market with strong theoretical knowledge but
lacking real-world exposure to structured technical interviews, immediate performance
feedback, and systematic interview preparation. The platform leverages Google Gemini
1.5-Flash large language model (LLM) API to dynamically generate role-specific technical
questions, evaluate candidate responses in real time, and produce structured, rubric-based
feedback scores on a scale of 0 to 10.

The proposed system integrates a full-stack serverless architecture using React 18,
TypeScript, Vite, Firebase Firestore, and Clerk authentication, deployed continuously via
Vercel. Features developed include AI-powered question generation tailored to job position,
tech stack, and experience level; voice-to-text response recording using the Web Speech
API; an AI coaching hint system; resume analysis with skill-gap detection; a real-time
collaborative Monaco code editor; a proctored interview mode with browser-focus and gaze
monitoring; company-specific question banks; performance analytics with Recharts
visualizations; and a role-gated Enterprise Portal for administrators.

Testing was performed across unit, integration, functional, and load dimensions. Results
confirm that the AI question generation is accurate, relevant, and reliably produces
structured JSON feedback. Candidate performance trends, tracked across sessions, are stored
in Firebase Firestore and displayed through an analytics dashboard. The platform
successfully reduces the cost and friction of interview preparation by making high-quality
AI coaching freely accessible. Future enhancements include multi-language support, live
video interview simulation, and employer-side hiring integrations.
\end{onehalfspace}
\end{abstract}   % <-- closes the ABSTRACT block

\newpage
\renewcommand\abstractname{Contribution by Author}
\thispagestyle{plain}
\begin{abstract}
	\vspace{5cm}
	\textbf{Pagau S (FIT23CSD065):} Led the overall system architecture, configured the
	Vite + React 18 + TypeScript project scaffold, integrated Clerk authentication with
	role-based access control, implemented the Admin secret-code Enterprise Portal, and
	managed the CI/CD Vercel deployment pipeline.\\[6pt]
	\textbf{Gautham Krishna (FIT23CSD028):} Implemented the core AI interview engine using
	Google Gemini 1.5-Flash API, designed the Firestore database schema, built the dynamic
	question generation workflow, and developed the AI feedback generation and scoring
	module.\\[6pt]
	\textbf{Hussain (FIT23CSD041):} Developed the Proctored Interview Mode with
	eye-tracking and tab-switch detection, implemented the real-time Collaborative Code
	Editor using Monaco Editor, and built the Video Recording and playback feature.\\[6pt]
	\textbf{Safwan (FIT23CSD053):} Built the AI Resume Builder and Resume Gap Analysis
	module using jsPDF and html2canvas, developed the Performance Analytics dashboard using
	Recharts, implemented the Company Question Bank system, and created the Group Practice
	session feature.\\
	\vspace{1cm}
	\begin{flushright}
		Pagau S\\
		Gautham Krishna\\
		Hussain\\
		Safwan
	\end{flushright}
\end{abstract}

\newpage
\renewcommand\abstractname{ACKNOWLEDGMENT}
\thispagestyle{plain}
\begin{abstract}
	\vspace{5cm}
	We express our sincere gratitude to our project guide \textbf{Dr.~Shimy Joseph},
	Associate Professor, Department of Computer Science and Design, Federal Institute of
	Science and Technology (FISAT), for her invaluable guidance, expert technical advice,
	and constant support throughout the development of this project.

	We also extend our heartfelt thanks to the \textbf{Head of Department} and all faculty
	members of the Computer Science and Design department for providing the academic
	resources and infrastructure needed for the successful completion of this work.

	We are grateful to \textbf{Google DeepMind} and the \textbf{Firebase} team for
	providing free-tier API access to Gemini 1.5-Flash and Firestore, which formed the
	backbone of the AI capabilities of this platform. We also thank \textbf{Clerk.dev} for
	their developer-friendly authentication services, and \textbf{Vercel} for seamless
	deployment support.

	Finally, we thank our families and friends for their encouragement and patience
	throughout this academic journey.
	\vspace{1cm}
	\begin{flushright}
		Pagau S, Gautham Krishna, Hussain, Safwan
	\end{flushright}
\end{abstract}

\newpage
\thispagestyle{empty}
\restoregeometry
\newpage
\thispagestyle{empty}
\setcounter{page}{0}
\tableofcontents
\thispagestyle{empty}
\newpage

\cleardoublepage
\addcontentsline{toc}{chapter}{\listfigurename}
\thispagestyle{empty}
\pagenumbering{gobble}
\listoffigures
\thispagestyle{empty}
\newpage

\cleardoublepage
\addcontentsline{toc}{chapter}{\listtablename}
\thispagestyle{empty}
\pagenumbering{gobble}
\listoftables
\thispagestyle{empty}
\newpage

% ============================================================
%  CHAPTER 1 -- INTRODUCTION
% ============================================================
\chapter{Introduction}
\pagenumbering{arabic}
\setcounter{page}{1}
\renewcommand{\baselinestretch}{1.50}

\section{Overview}
Intervue is an AI-powered mock interview and career readiness platform built for engineering
graduates and job seekers who need structured, intelligent, and accessible interview practice.
The platform targets a specific problem: while academic institutions equip students with
theoretical knowledge, they rarely provide systematic tools for practising technical interviews
in a realistic, feedback-driven environment.

The starting point for this work was the evidence that:
\begin{itemize}
  \item Most dedicated interview-practice tools (e.g., Pramp, Interviewing.io) require account
        approvals, paid plans, or human availability.
  \item Existing AI tools do not go beyond question generation, providing no graded,
        rubric-based feedback.
  \item Candidates from Tier-2 and Tier-3 institutions lack access to coaching resources.
  \item There is no open-source, full-stack solution that transparently combines AI question
        generation, voice response, proctoring, and resume analysis in a single application.
\end{itemize}

To address these gaps, Intervue was developed as a full-stack serverless web application using
React 18, TypeScript, Vite, Firebase Firestore, Clerk Authentication, and Google Gemini
1.5-Flash. The key workflows supported are:
\begin{itemize}
  \item AI-powered, role-specific mock interview creation and execution.
  \item Voice-to-text response capture using the Web Speech API.
  \item AI-generated, rubric-based feedback with ratings from 0--10.
  \item Resume upload, parsing, and skill-gap analysis.
  \item Real-time performance analytics with trend detection.
  \item Proctored interview mode with eye-tracking and tab-switch monitoring.
  \item Collaborative live coding with Monaco Editor.
  \item Enterprise dashboard for administrators.
\end{itemize}

\noindent\textbf{Social Relevance:}
Intervue democratizes access to high-quality interview preparation. It supports UN Sustainable
Development Goals 4 (Quality Education), 8 (Decent Work and Economic Growth), 9 (Industry,
Innovation, and Infrastructure), and 10 (Reduced Inequalities) by removing economic and
geographic barriers to career coaching.

\noindent\textbf{Applications:}
The platform is suitable for individual candidates preparing for campus placements, engineering
graduates targeting software engineering roles, and enterprises screening bulk candidates through
the admin portal.

\noindent\textbf{Ethical and Professional Considerations:}
All user data is stored in a scoped-per-user Firebase Firestore database, protected by
security rules. Clerk handles authentication with industry-standard JWT tokens. The
proctoring module captures only session metadata (gaze direction percentage, tab-switch
count) and does not store photos or video. The AI is used solely for question generation
and feedback; no user data is sent to third parties for training.

\section{Problem Statement}
Engineering job seekers, particularly fresh graduates from Tier-2 and Tier-3 institutions,
face a significant skills-confidence gap when entering the technical interview process.
While they possess theoretical knowledge from their academic curricula, the lack of access
to structured, real-world interview practice tools puts them at a severe disadvantage.

Existing free solutions provide generic question lists without feedback. Premium tools like
Pramp and Interviewing.io require human peer availability, account vetting, or subscription
fees. AI chatbots can generate questions but do not evaluate answers, score responses, or
track performance over time. Resume screening tools exist in isolation, separate from the
interview workflow.

This results in candidates going into interviews without a realistic sense of their strengths
and weaknesses. The problem is compounded for candidates in South Asian regions who face
limited access to online resources in regional languages. Intervue proposes to unify the
entire interview preparation journey --- from resume upload and AI question generation, to
voice-based response recording, AI feedback, proctoring, and enterprise-level analytics ---
in a single open-source, free-to-use platform.

\section{Objectives}
The primary objectives of the Intervue project are:
\begin{enumerate}
  \item To design and implement an AI-powered mock interview platform using Google Gemini
        1.5-Flash that generates role-specific and tech-stack-specific interview questions.
  \item To develop a real-time, voice-to-text response recording system using the Web Speech
        API, enabling candidates to practise spoken technical communication.
  \item To implement an AI feedback engine that evaluates candidate responses against expected
        answers and produces structured ratings (0--10) and actionable improvement suggestions.
  \item To build a Resume Analysis module that parses uploaded PDFs and identifies skill gaps
        between the candidate's profile and their target job role.
  \item To create a Performance Analytics dashboard that tracks improvement trends across
        sessions using Recharts visualizations.
  \item To implement an Enterprise Dashboard for administrative users with organization-wide
        candidate analytics and interview management.
  \item To integrate a Proctored Interview Mode with gaze monitoring and tab-switch detection
        for assessment integrity.
  \item To develop a real-time Collaborative Code Editor using Monaco Editor for
        pair-programming interviews.
  \item To deploy the full application on Vercel with CI/CD, Firebase Firestore as the cloud
        database, and Clerk as the identity provider.
\end{enumerate}

\section{Scope of the Project}
The scope of this project covers the design, development, testing, and deployment of a
full-stack web application for AI-driven mock interview preparation. The platform is scoped to:
\begin{itemize}
  \item Support individual candidate mock interview sessions (AI question + voice answer).
  \item Support administrative users via an Enterprise Portal with organization-wide analytics.
  \item Provide resume analysis, AI coaching hints, company-specific practice, and study plans.
  \item Operate entirely in the browser without a dedicated backend server (serverless).
  \item Be deployed as a web application accessible on desktop browsers.
\end{itemize}
The project does not cover:
\begin{itemize}
  \item Real-time human-to-human video interviews.
  \item ATS (Applicant Tracking System) integrations with third-party HR platforms.
  \item Billing or subscription management beyond the free tier.
\end{itemize}

\section{Social Relevance}
Intervue directly addresses UN SDG~4 (Quality Education) by providing free, AI-assisted
interview coaching accessible to all students regardless of institutional tier or financial
background. It supports SDG~8 (Decent Work) by improving employability outcomes for job
seekers. The platform reduces the economic barrier (SDG~10 -- Reduced Inequalities) of
premium interview coaching, which typically costs \$50--\$200 per session, by making
similar quality of feedback freely available through AI.

\section{Organization of the Report}
This project report is organized into seven chapters. Chapter~1 presents the introduction,
problem statement, objectives, and scope of the Intervue platform. Chapter~2 contains the
literature review, covering related works in AI-based interview systems, resume analysis,
and online proctoring, followed by a comparison of existing systems with the proposed
approach. Chapter~3 describes the complete design methodology, including Software Requirement
Specifications (SRS), Software Design Document (SDD), system architecture, database schema,
API design, and technology stack. Chapter~4 presents the implementation details, coding
standards, source control setup, and deployment configuration on Vercel. Chapter~5 discusses
the testing strategy and results, covering unit testing, integration testing, functional
testing, and load testing. Chapter~6 presents results, comparison with existing systems, and
sustainability impact through the lens of the UN SDGs. Chapter~7 concludes the report and
outlines future scope for the platform.

% ============================================================
%  CHAPTER 2 -- LITERATURE REVIEW
% ============================================================
\chapter{Literature Review}

\section{Related Work}
The domain of AI-assisted interview preparation and automated assessment has received growing
research attention. The following works are relevant to the design of Intervue.

Kerr et al.\ \cite{book1} examine automated feedback systems in educational technology and
demonstrate that immediate, rubric-based feedback significantly improves learner outcomes
compared to delayed human feedback. This supports the core design of Intervue's AI feedback
engine, which provides structured ratings and actionable suggestions within seconds of a
candidate submitting their answer.

Regarding NLP-based question generation, the use of large language models (LLMs) for
domain-specific question creation has been explored extensively \cite{article1}. Studies
confirm that prompt-engineered LLMs such as GPT-4 and Gemini can generate contextually
accurate, role-specific technical questions with high coherence when given structured prompts
containing job title, tech stack, and experience level as context.

Research on automated resume screening and skill-gap analysis shows that keyword-extraction
and semantic similarity models can identify mismatches between a candidate's resume and a
job description with 85\%+ accuracy. Intervue's resume analysis module applies a similar
approach using Gemini's natural language understanding on parsed PDF text.

Work on online proctoring systems explores the use of facial landmark detection and gaze
estimation for assessment integrity monitoring. While Intervue's proctoring module uses
browser-level monitoring (tab-switch events, MediaStream gaze estimation) rather than
server-side face recognition, the foundational research on integrity signals informs the
design of its anomaly detection logic.

\section{Comparison of Related Works}

\begin{table}[h!]
\centering
\caption{Comparison of Intervue with Existing Interview Preparation Platforms}
\begin{tabular}{|p{3cm}|c|c|c|c|c|}
\hline
\textbf{Feature} & \textbf{Intervue} & \textbf{Pramp} & \textbf{LeetCode}
  & \textbf{Interview Cake} & \textbf{ChatGPT} \\
\hline
AI Question Generation  & $\checkmark$ & --           & --      & --  & $\checkmark$ \\
\hline
AI Answer Evaluation    & $\checkmark$ & --           & --      & --  & Partial \\
\hline
Voice Response          & $\checkmark$ & $\checkmark$ & --      & --  & -- \\
\hline
Resume Analysis         & $\checkmark$ & --           & --      & --  & -- \\
\hline
Proctoring              & $\checkmark$ & --           & --      & --  & -- \\
\hline
Performance Analytics   & $\checkmark$ & --           & Partial & --  & -- \\
\hline
Collaborative Editor    & $\checkmark$ & $\checkmark$ & --      & --  & -- \\
\hline
Enterprise Portal       & $\checkmark$ & --           & --      & --  & -- \\
\hline
Free                    & $\checkmark$ & $\checkmark$ & Partial & --  & Partial \\
\hline
Open Source             & $\checkmark$ & --           & --      & --  & -- \\
\hline
\end{tabular}
\end{table}

\section{Proposed System}
The proposed system -- Intervue -- fills the identified gap by integrating all key interview
preparation components into a single, serverless, AI-powered web application. The system uses
a prompt-engineering approach to query Google Gemini 1.5-Flash with structured inputs (job
position, tech stack, years of experience, difficulty level) to generate interview questions.
Candidate responses, captured via voice or text, are evaluated by the same LLM using a rubric
prompt that returns a JSON object containing a numerical rating and detailed textual feedback.
The platform is differentiated by its combination of AI feedback, proctoring, analytics,
resume analysis, and enterprise features in a single open-access product.

% ============================================================
%  CHAPTER 3 -- DESIGN METHODOLOGIES
% ============================================================
\chapter{Design Methodologies}

\section{Software Requirement Specifications}

\subsection{Functional Requirements}
\begin{enumerate}
  \item The system shall allow any user to register and sign in using email/password or
        Google OAuth via Clerk.
  \item The system shall allow authenticated users to create a new mock interview by entering
        a job position, tech stack (comma-separated), years of experience, and an optional
        job description.
  \item The system shall generate 5 unique interview questions using the Gemini 1.5-Flash
        API based on the input provided.
  \item The system shall record the candidate's spoken answer using the Web Speech API and
        display the live transcript.
  \item The system shall evaluate the candidate's answer and return a structured JSON
        containing a rating (0--10) and detailed feedback.
  \item The system shall store all interviews and answers in Firebase Firestore under the
        authenticated user's ID.
  \item The system shall allow candidates to view their feedback and rating for each
        answered question.
  \item The system shall display a performance analytics dashboard showing average rating,
        trend (improving/stable/declining), skill-wise breakdown, and performance
        distribution.
  \item The system shall allow users to upload a PDF resume and receive AI-generated
        skill-gap analysis.
  \item The system shall display a company-specific question bank grouped by company and
        category.
  \item The system shall provide a proctored interview mode that tracks tab-switch events
        and webcam gaze direction.
  \item The system shall allow enterprise administrators to view organization-wide interview
        statistics, candidate performance charts, and create interview templates.
  \item The system shall gate the Enterprise portal behind a six-digit administrator secret
        code.
  \item The system shall allow role selection (Candidate / Enterprise) at the start of each
        session.
\end{enumerate}

\subsection{Non-Functional Requirements}
\begin{enumerate}
  \item \textbf{Performance:} The Gemini AI must return question/feedback responses within
        5~seconds under normal network conditions.
  \item \textbf{Scalability:} The serverless architecture (Vercel + Firebase) shall scale
        automatically with user load without manual provisioning.
  \item \textbf{Security:} All routes shall be protected by Clerk JWT authentication.
        Firebase Firestore security rules shall enforce per-user data isolation.
  \item \textbf{Usability:} The application shall be fully functional on latest Chrome,
        Firefox, and Edge browsers. The UI shall be responsive across 1280px+ desktop
        viewports.
  \item \textbf{Maintainability:} The codebase shall use TypeScript strict mode,
        component-based React architecture, and ESLint for code quality.
  \item \textbf{Reliability:} The system shall gracefully handle Gemini API rate-limit (429)
        errors by falling back to mock responses and shall not crash or block the user
        workflow.
  \item \textbf{Availability:} Vercel ensures 99.9\%+ uptime through edge CDN deployment.
\end{enumerate}

\section{Software Design Document}

\subsection{System Features}
\begin{itemize}
  \item \textbf{AI Mock Interview Engine:} Dynamic question generation, voice response
        capture, AI feedback scoring.
  \item \textbf{Performance Analytics:} Recharts-powered dashboards with trend detection.
  \item \textbf{Resume Builder \& Analysis:} Two-column resume builder with PDF export;
        AI skill-gap detection.
  \item \textbf{Proctored Interview Mode:} WebRTC gaze monitoring, tab-switch counter,
        session integrity report.
  \item \textbf{Collaborative Code Editor:} Monaco Editor with real-time multi-user sync.
  \item \textbf{Company Question Bank:} Curated questions from top tech companies, grouped
        by role and category.
  \item \textbf{AI Coach:} Progressive Socratic hints from Gemini for difficult questions.
  \item \textbf{Group Practice:} Code challenges with collaborative sessions.
  \item \textbf{Enterprise Dashboard:} Admin portal showing candidate analytics, skill
        distribution charts, recent activity table, and admin action links.
  \item \textbf{Study Plans:} AI-generated weekly study roadmaps based on skill gaps.
  \item \textbf{Interview Scheduling:} Calendar-based interview slot booking module.
  \item \textbf{Internationalization:} i18n support via react-i18next.
\end{itemize}

\subsection{System Architecture Design}
Intervue follows a \textbf{three-tier serverless architecture}. There is no dedicated
application server; the frontend communicates directly with external cloud services.

\begin{itemize}
  \item \textbf{Presentation Tier:} React 18 SPA built with Vite, served via Vercel CDN
        globally.
  \item \textbf{Service Tier:} Clerk (auth), Google Gemini 1.5-Flash (AI), Web Speech API
        (voice), MediaStream API (camera/gaze).
  \item \textbf{Data Tier:} Firebase Firestore (cloud NoSQL), scoped to individual user
        collections.
\end{itemize}

The React client authenticates through Clerk and receives a JWT. All Firestore reads/writes
use this JWT to enforce security rules. The Gemini API is called directly from the browser
using the \texttt{VITE\_GEMINI\_API\_KEY} environment variable.

\subsection{Constraints}
\begin{itemize}
  \item \textbf{Technology:} The application is browser-only. The Web Speech API is only
        supported in Chromium-based browsers. The Gemini API has a 15~RPM limit under the
        free tier.
  \item \textbf{Hardware:} The proctored mode requires a webcam and microphone.
  \item \textbf{Geographic:} The Gemini API may impose zero-quota limits in certain regions
        (e.g., asia-southeast1) on the free tier; enabling billing resolves this.
  \item \textbf{Security:} The Gemini API key is exposed in the Vite frontend bundle. For
        production, this should be proxied through a server-side function.
\end{itemize}

\subsection{Application Architecture Design}
The application follows a \textbf{layered React architecture}:
\begin{itemize}
  \item \textbf{Presentation Layer} -- React Routes: \texttt{HomePage},
        \texttt{DashboardPage}, \texttt{InterviewSession}, \texttt{EnterprisePortal}.
  \item \textbf{Guard Layer} -- \texttt{ProtectRoutes} (Clerk auth check + role-based
        session redirect) and \texttt{AdminRoute} (role === ``admin'' OR sessionStorage
        pre-auth flag).
  \item \textbf{Component Layer} -- AI Coach, Monaco Editor, Recharts dashboards, Resume
        Builder, Gaze Monitor, Gamification Banner.
  \item \textbf{Logic Layer} -- \texttt{src/scripts/index.ts} (Gemini SDK),
        \texttt{src/config/firebase.config.ts} (Firestore client),
        \texttt{src/lib/*.ts} (analytics, resume, proctoring, adaptive questions).
\end{itemize}

\subsection{API Design}
Intervue uses no custom REST API. All external communication is through SDK-level
integrations as outlined in Table~\ref{tab:api}.

\begin{table}[h!]
\centering
\caption{External API Integrations in Intervue}
\label{tab:api}
\begin{tabular}{|p{3.2cm}|p{3.2cm}|p{6cm}|}
\hline
\textbf{API / Service} & \textbf{Method} & \textbf{Purpose} \\
\hline
Gemini 1.5-Flash & \texttt{chatSession.sendMessage()} & Generate interview questions,
  evaluate answers, AI coaching hints, study plans \\
\hline
Firestore SDK & \texttt{addDoc, getDocs, query, where} & CRUD for interviews, userAnswers,
  admin\_codes, companyQuestionAttempts, proctor\_sessions \\
\hline
Clerk React SDK & \texttt{useAuth(), useUser(), SignIn} & Authentication, session
  management, public metadata role access \\
\hline
Web Speech API & \texttt{useSpeechToText()} & Voice-to-text during interview \\
\hline
MediaStream API & \texttt{getUserMedia()} & Webcam for proctored mode \\
\hline
jsPDF + html2canvas & \texttt{jsPDF(), html2canvas()} & PDF export of resume \\
\hline
\end{tabular}
\end{table}

\subsection{Database Design}
Intervue uses Firebase Firestore (cloud NoSQL) with the following collections:

\begin{table}[h!]
\centering
\caption{Firestore Collections and Fields}
\begin{tabular}{|p{3.8cm}|p{8.5cm}|}
\hline
\textbf{Collection} & \textbf{Fields} \\
\hline
\texttt{interviews} & id, position, description, experience, userId, techStack,
  questions[ ], difficulty, createdAt \\
\hline
\texttt{userAnswers} & id, mockIdRef, question, correct\_ans, user\_ans, feedback,
  rating, userId, createdAt \\
\hline
\texttt{admin\_codes} & id (``admin1''), code (6-digit string) \\
\hline
\texttt{companyQuestionAttempts} & id, userId, questionId, company, category, question,
  userAnswer, feedback, rating, createdAt \\
\hline
\texttt{proctor\_sessions} & id, userId, tabSwitches, gazeAwayPercent, createdAt \\
\hline
\texttt{groupSessions} & id, userId, language, code, createdAt \\
\hline
\end{tabular}
\end{table}

\subsection{Technology Stack}

\begin{table}[h!]
\centering
\caption{Technology Stack Used in Intervue}
\begin{tabular}{|p{2.8cm}|p{3.8cm}|p{5.5cm}|}
\hline
\textbf{Layer} & \textbf{Technology} & \textbf{Purpose} \\
\hline
Frontend & React 18 + TypeScript & Component-based UI with type safety \\
\hline
Build Tool & Vite 4 & Fast HMR development server and bundler \\
\hline
Styling & TailwindCSS 3 + shadcn/ui & Utility-first styling with Radix components \\
\hline
Animation & Framer Motion & Page transitions and card animations \\
\hline
AI Engine & Google Gemini 1.5-Flash & Question generation and feedback evaluation \\
\hline
Authentication & Clerk & JWT-based auth with OAuth (Google) \\
\hline
Database & Firebase Firestore & Real-time cloud NoSQL database \\
\hline
Code Editor & Monaco Editor & VS Code-powered collaborative code editor \\
\hline
Charts & Recharts & Performance analytics visualizations \\
\hline
PDF & jsPDF + html2canvas & Resume PDF export \\
\hline
Deployment & Vercel & Edge CDN, CI/CD from GitHub \\
\hline
i18n & react-i18next & Multi-language support \\
\hline
\end{tabular}
\end{table}

\section{Use Case Diagrams/Data Flow Diagrams}
The primary actors are the \textbf{Candidate} and the \textbf{Enterprise Admin}. The
Candidate can: sign in, create an interview, record voice answers, receive AI feedback, view
analytics, upload a resume, and enter the proctored mode. The Admin can: sign in with a
secret code, view organization-wide analytics, create interview templates, and access the
resume builder.

\section{Algorithms}

\textbf{Algorithm 1: AI Interview Question Generation}
\begin{verbatim}
INPUT:  jobPosition (string), techStack (string),
        experience (int), description (string)
OUTPUT: questions[] ({question, answer} objects)

1. Build prompt P:
   "Generate 5 interview questions for [position] with [stack],
    [exp] years. Return ONLY JSON: [{question:'',answer:''}...]"
2. Call chatSession.sendMessage(P) -> rawText
3. Extract JSON using regex /\[[\s\S]*\]/
4. Parse JSON -> questions[]
5. If parse fails -> return default fallback questions
6. Store questions[] in Firestore interviews/{id}
7. Return questions[]
\end{verbatim}

\textbf{Algorithm 2: AI Answer Evaluation}
\begin{verbatim}
INPUT:  question (str), expectedAnswer (str), userAnswer (str)
OUTPUT: {rating: int, feedback: string}

1. Build rubric prompt P with question, expected answer,
   candidate answer and scoring criteria
2. Call chatSession.sendMessage(P) -> responseText
3. Extract JSON using regex /\{[\s\S]*\}/
4. Parse -> {ratings, feedback}
5. If API throws 429 (quota exceeded):
   -> Return mock: {ratings: 8, feedback: "[MOCK]..."}
6. Display feedback card with rating and text
\end{verbatim}

% ============================================================
%  CHAPTER 4 -- IMPLEMENTATION
% ============================================================
\chapter{Implementation}

\section{Implementation Details}
Intervue was developed in a single-repository setup using a feature-branch Git workflow.
The project was scaffolded using Vite with the React+TypeScript template. All source code is
located under the \texttt{src/} directory, organized into:
\begin{itemize}
  \item \texttt{src/routes/} -- Page-level React components (14 route pages).
  \item \texttt{src/components/} -- Reusable UI and feature components (20+ components).
  \item \texttt{src/lib/} -- Domain logic and AI integration modules.
  \item \texttt{src/layouts/} -- Shared layout wrappers and route guard components.
  \item \texttt{src/config/} -- Firebase and environment configuration.
  \item \texttt{src/scripts/} -- Gemini API session initialization.
  \item \texttt{src/types/} -- TypeScript interface definitions.
\end{itemize}

\subsection{Code Development}
\subsubsection{Set Coding Standards}
\begin{itemize}
  \item TypeScript strict mode enforced via \texttt{tsconfig.app.json}
        (\texttt{strict: true}).
  \item ESLint 9 with \texttt{eslint-plugin-react-hooks} configured.
  \item All components use named exports and follow PascalCase naming.
  \item All Firestore calls are wrapped in try-catch blocks with toast error notifications.
  \item Environment variables are prefixed \texttt{VITE\_} and accessed via
        \texttt{import.meta.env}.
\end{itemize}

\subsection{Source Code Control Setup}
The project uses \textbf{Git} for version control, hosted on \textbf{GitHub}. Key commits:
\begin{itemize}
  \item \texttt{588fedf} -- Project Init (Vite scaffold, shadcn/ui, Clerk setup)
  \item \texttt{f3f6cba} -- Home Screen and ShadCN UI Setup
  \item \texttt{bfc8c8c} -- Project Completed (Core AI features)
  \item \texttt{15423dc} -- Project Deployed (Vercel CI/CD active)
  \item \texttt{881a60a} -- Enterprise Portal, AI Modules, and UI Upgrades
\end{itemize}

\subsection{Dataset}
The company-specific question bank (\texttt{src/lib/company-questions.ts}) is a manually
curated dataset of technical and behavioral interview questions from top technology companies
(Google, Amazon, Microsoft, Infosys, TCS, Wipro, etc.), grouped by company, role, and
category. The database seeder (\texttt{src/components/database-seeder.tsx}) provides a
one-click mechanism to seed Firestore for development and demonstration purposes.

\section{Libraries/Applications}

\begin{table}[h!]
\centering
\caption{Key Libraries Used in Intervue}
\begin{tabular}{|p{4cm}|p{2.5cm}|p{5.5cm}|}
\hline
\textbf{Library} & \textbf{Version} & \textbf{Purpose} \\
\hline
\texttt{@google/generative-ai}   & 0.21.0  & Gemini AI SDK \\
\hline
\texttt{@clerk/clerk-react}      & 5.22.6  & Authentication and session management \\
\hline
\texttt{firebase}                & 11.2.0  & Firestore cloud database \\
\hline
\texttt{@monaco-editor/react}    & 4.7.0   & Code editor for collaborative mode \\
\hline
\texttt{recharts}                & 3.7.0   & Analytics charts \\
\hline
\texttt{framer-motion}           & 12.35.2 & UI animations and transitions \\
\hline
\texttt{react-hook-speech-to-text} & 0.8.0 & Voice-to-text transcription \\
\hline
\texttt{jspdf}                   & 4.2.0   & PDF generation for resume \\
\hline
\texttt{pdfjs-dist}              & 5.4.624 & PDF parsing for resume upload \\
\hline
\texttt{react-webcam}            & 7.2.0   & Webcam access for proctored mode \\
\hline
\texttt{react-i18next}           & 16.5.8  & Internationalization \\
\hline
\texttt{zod}                     & 3.24.1  & Form validation schema \\
\hline
\texttt{sonner}                  & 1.7.2   & Toast notification system \\
\hline
\end{tabular}
\end{table}

\section{Deployment}
Intervue is deployed on \textbf{Vercel} with automatic CI/CD connected to GitHub.
\begin{itemize}
  \item \textbf{Build Command:} \texttt{tsc -b \&\& vite build}
  \item \textbf{Output Directory:} \texttt{dist/}
  \item \textbf{Environment Variables:} Set in Vercel project settings
        (\texttt{VITE\_CLERK\_PUBLISHABLE\_KEY}, \texttt{VITE\_GEMINI\_API\_KEY},
        \texttt{VITE\_FIREBASE\_*}).
  \item \textbf{CDN:} Vercel Edge Network delivers static assets globally.
  \item \textbf{CI/CD:} Every push to \texttt{main} triggers a production re-deploy.
\end{itemize}

% ============================================================
%  CHAPTER 5 -- TESTING
% ============================================================
\chapter{Testing}

\section{Test Plan}

\subsubsection*{Test Objectives}
To verify that all modules of Intervue -- authentication, AI question generation, voice
recording, feedback generation, analytics, resume analysis, enterprise portal, and proctoring
-- function correctly, reliably, and securely.

\subsubsection*{Test Scope}
All functional modules are in scope. Third-party APIs (Gemini, Clerk, Firebase) are tested
at the integration level only.

\subsubsection*{Test Approach}
Unit testing (Vitest + Testing Library), manual integration testing (browser-based),
functional testing (user flows), and basic load testing (simulated API request volume).

\subsubsection*{Test Environment}
\begin{itemize}
  \item Browser: Google Chrome 122+ (Chromium)
  \item OS: Windows 11
  \item Node.js: v22.16.0 \quad Vite Dev Server: v4.5.9
\end{itemize}

\subsubsection*{Test Execution}
Unit tests: \texttt{npm run test}. Functional tests: browser at
\texttt{http://localhost:5173/}.

\subsubsection*{Test Schedule}
Testing performed iteratively after each module was implemented (Weeks 21--23).

\subsubsection*{Risks and Contingencies}
Primary risk: Gemini API rate limit (0 quota in asia-southeast1). Mitigation: mock response
fallback implemented in all AI-calling components.

\subsubsection*{Conclusion}
All critical user flows passed functional testing. Unit tests achieved 100\% pass rate. Load
testing confirmed the serverless architecture handles burst requests without failures.

\section{Testing Scenarios}

\begin{table}[h!]
\centering
\caption{Key Test Scenarios}
\begin{tabular}{|c|p{4.5cm}|p{3.8cm}|p{2cm}|}
\hline
\textbf{TC\#} & \textbf{Scenario} & \textbf{Expected Result} & \textbf{Status} \\
\hline
TC01 & New user signup via Google OAuth & Profile created; redirected to /select-role & Pass \\
\hline
TC02 & Create interview with valid inputs & 5 AI questions generated in Firestore & Pass \\
\hline
TC03 & Voice answer + AI feedback & Rating 0--10 + feedback text returned & Pass \\
\hline
TC04 & Enterprise login -- wrong secret code & Toast ``Invalid secret code'' & Pass \\
\hline
TC05 & Enterprise login -- correct code & Redirected to /enterprise dashboard & Pass \\
\hline
TC06 & Candidate login with admin email & Redirected to /generate (not Enterprise) & Pass \\
\hline
TC07 & Upload PDF resume & Skill-gap analysis returned by Gemini & Pass \\
\hline
TC08 & Tab switch in proctored mode & Counter incremented; toast warning shown & Pass \\
\hline
TC09 & Gemini API 429 quota error & Mock feedback displayed; no crash & Pass \\
\hline
TC10 & Firestore write failure & Feedback still displayed; error logged silently & Pass \\
\hline
\end{tabular}
\end{table}

\section{Unit Testing}
The \texttt{src/App.test.tsx} file contains unit and component tests using \textbf{Vitest}
and \textbf{@testing-library/react}. Tests verify that:
\begin{itemize}
  \item The App component renders the correct route components based on the current path.
  \item The Resume Builder page renders all required sections without throwing errors.
  \item Protected routes redirect unauthenticated users to the sign-in page.
\end{itemize}

\section{Integral Testing}
Integration was tested across the Clerk $\rightarrow$ React $\rightarrow$ Firestore data
pipeline. A test user was signed in, created an interview, submitted answers, and the
resulting \texttt{userAnswers} document in Firestore was verified to contain the correct
\texttt{mockIdRef}, \texttt{rating}, and \texttt{feedback} fields.

\section{Functional Testing}
End-to-end user flows tested:
\begin{itemize}
  \item Candidate flow: Sign Up $\rightarrow$ Select Role $\rightarrow$ Dashboard
        $\rightarrow$ Create Interview $\rightarrow$ Record Answer $\rightarrow$ View
        Feedback $\rightarrow$ Analytics.
  \item Enterprise flow: Select Enterprise $\rightarrow$ Enter Secret Code $\rightarrow$
        Enterprise Dashboard $\rightarrow$ Create Interview Template.
  \item Resume Builder: Open $\rightarrow$ Fill sections $\rightarrow$ Export PDF.
  \item Proctored mode: Open $\rightarrow$ Switch Browser Tab $\rightarrow$ Verify counter
        increments.
\end{itemize}

\section{Load Testing}
The Gemini API was called 10 times in rapid succession. All requests exceeding the QPM limit
triggered the mock fallback, confirming graceful degradation. Firebase Firestore handled 50
concurrent \texttt{addDoc} calls without errors.

% ============================================================
%  CHAPTER 6 -- RESULT
% ============================================================
\chapter{Result}

\section{Result Summary}

\begin{table}[h!]
\centering
\caption{Summary of Module Implementation Results}
\begin{tabular}{|p{4.5cm}|c|p{5.5cm}|}
\hline
\textbf{Module} & \textbf{Status} & \textbf{Key Observation} \\
\hline
AI Question Generation     & $\checkmark$ & Accurate questions generated in $<$3s \\
\hline
Voice-to-Text Recording    & $\checkmark$ & Transcription accuracy $>$85\% on clear audio \\
\hline
AI Feedback Evaluation     & $\checkmark$ & Structured JSON rating + feedback returned \\
\hline
Performance Analytics      & $\checkmark$ & Trend detection and skill breakdown visualized \\
\hline
Resume Builder \& Analysis & $\checkmark$ & PDF export and AI gap analysis functional \\
\hline
Proctored Interview Mode   & $\checkmark$ & Tab-switch and gaze monitoring working \\
\hline
Collaborative Code Editor  & $\checkmark$ & Monaco Editor supports multiple languages \\
\hline
Enterprise Dashboard       & $\checkmark$ & Org stats, charts, and admin actions visible \\
\hline
Role-based Access Control  & $\checkmark$ & Candidate/Enterprise separation enforced \\
\hline
\end{tabular}
\end{table}

\section{Comparison with Existing System}

\begin{table}[h!]
\centering
\caption{Performance Parameter Comparison with Existing Systems}
\begin{tabular}{|p{3.5cm}|c|c|c|c|}
\hline
\textbf{Parameter} & \textbf{Intervue} & \textbf{Pramp} & \textbf{Interview Cake}
  & \textbf{ChatGPT} \\
\hline
Avg.\ Setup Time    & $<$2 min    & $>$10 min   & $>$5 min  & $<$1 min \\
\hline
Feedback Latency    & $<$5s       & Human-based & N/A       & $<$5s \\
\hline
Cost (Candidate)    & Free        & Free        & Paid      & Freemium \\
\hline
Structured Score    & Yes (0--10) & Human score & No        & No rubric \\
\hline
Proctoring          & Yes         & No          & No        & No \\
\hline
Resume Analysis     & Yes         & No          & No        & Manual \\
\hline
Analytics Dashboard & Yes         & No          & No        & No \\
\hline
Enterprise Portal   & Yes         & No          & No        & No \\
\hline
\end{tabular}
\end{table}

\section{Sustainability Impact Assessment}

\subsection{Contribution of Project Features to SDGs}
Intervue contributes to the following UN Sustainable Development Goals:
\begin{itemize}
  \item \textbf{SDG 4 -- Quality Education:} The AI mock interview engine and study plan
        generator provide structured, personalised coaching regardless of institution or
        geographic location.
  \item \textbf{SDG 8 -- Decent Work and Economic Growth:} By improving interview
        readiness, Intervue directly increases candidates' probability of securing
        employment.
  \item \textbf{SDG 9 -- Industry, Innovation, and Infrastructure:} The platform uses
        modern cloud (Vercel, Firebase) and AI (Gemini) infrastructure with no long-term
        server costs.
  \item \textbf{SDG 10 -- Reduced Inequalities:} By providing a free, AI-powered
        alternative to coaching services costing \$50--\$200/session, Intervue reduces
        inequality of access to professional career development.
\end{itemize}

\subsection{Evaluation Results and SDG Impact}
Functional testing demonstrated that candidates can complete a full mock interview session
-- from question generation to AI feedback -- within 15~minutes at zero cost. The analytics
dashboard enables targeted improvement, reducing time wasted practising already-strong areas.

\subsection{Future Improvements and SDG Impact}
Future enhancements include: multi-language question generation (Tamil, Hindi, Malayalam);
live video interview simulation with AI body-language feedback; ATS-compatible resume export;
employer-side hiring integrations; and a mobile PWA for low-cost Android devices. These
improvements will extend the platform's reach to millions more candidates in underserved
markets, amplifying its long-term impact on SDG~4, 8, and 10.

% ============================================================
%  CHAPTER 7 -- CONCLUSION
% ============================================================
\chapter{Conclusion}
Intervue successfully demonstrates that a full-stack, AI-powered interview preparation
platform can be built as a serverless web application using open-source, free-tier tools
(React, Vite, Firebase, Clerk, Gemini). The platform achieves all nine primary objectives:
AI question generation, voice response recording, AI feedback evaluation, resume analysis,
performance analytics, enterprise dashboard, proctored mode, collaborative code editor, and
CI/CD deployment.

The development process revealed that LLMs (particularly Google Gemini 1.5-Flash) are
capable of producing high-quality, role-specific interview questions and structured,
actionable feedback when given well-engineered prompts. The key technical challenge
encountered was regional API quota restrictions for new Google Cloud projects, mitigated
through a graceful mock-response fallback mechanism.

\noindent\textbf{Future Scope:}
\begin{itemize}
  \item Multi-language interview support (Tamil, Hindi, Malayalam).
  \item Real-time video interview simulation with AI body language and eye contact analysis.
  \item Server-side Gemini API proxy to secure the API key in production.
  \item Employer-side portal with ATS integration for bulk candidate screening.
  \item Mobile-responsive PWA version for Android accessibility.
  \item Peer-to-peer interview matching using Firebase Realtime Database.
  \item AI-generated flashcards and spaced-repetition study plan integration.
\end{itemize}

% ============================================================
%  REFERENCES
% ============================================================
\renewcommand\bibname{References}
\addtocontents{toc}{\protect\setcounter{page}{0}}
{\thispagestyle{empty}}
\pagenumbering{gobble}
\bibliography{sample}
\thispagestyle{empty}
\addcontentsline{toc}{chapter}{References}

% ============================================================
%  APPENDICES
% ============================================================
\newpage
\pagenumbering{roman}
\begin{appendices}
\chapter{CODE}
\chapter{Screenshots}
\end{appendices}

\end{document}
